\documentclass[UTF8,11pt]{ctexart}
\usepackage{amsmath,amssymb,geometry}
\geometry{margin=22mm}
\begin{document}
\section*{P1: 多站点患者预约的样本平均混合二元线性模型}
\paragraph{来源与范围}依据原论文PDF第1, 6, 7, 8, 9, 10, 11, 12页，(4)–(9)、Remark 1及命题2 (11)–(25)。
\subsection*{集合、参数与范围}
患者$k=1,\ldots,K$按固定相同顺序经过单服务者站点$n=1,\ldots,N$；所有患者均到诊且依次访问全部站点。$\Omega$是给定的有限等权服务时长情景样本，$\delta_k^n(\omega)\in\mathbb Z_+$是时隙数，允许跨患者和站点相关。$T^n\in\mathbb Z_+$为站点正常工时。费用满足$0\leq C_{L,k}\leq C_{W,k}$、$0\leq C_R^n\leq C_I^n$、$C_O^n\geq0$，依次对应休闲/普通等待、可利用/普通空闲和加班。可用的大常数是
\[
M=\max\{1,\max_nT^n+\max_{\omega\in\Omega}\sum_{k=1}^K\sum_{n=1}^N\delta_k^n(\omega)\}.
\]
本文件采用论文命题2的样本平均混合二元线性形式，预约时间$a_k^n$是连续变量（原文Remark 1允许放松整数时隙，命题2也直接使用连续域）。目标为给定服务时长样本上的平均费用。
\subsection*{变量与唯一完整模型}
$a$是事前预约表；情景变量$r$为患者就绪时刻，$v$为医护可接诊时刻，$s$为实际开始时刻，$f$为下班时刻。$y,z$为max关系的二元分支变量；$\Gamma_W,\Gamma_L,\Gamma_I,\Gamma_R,\Gamma_O$依次为普通等待、休闲等待、普通空闲、可利用空闲和加班的合计时长。边界处无效的$y,z$可不建立。
\begin{equation}
\min
\frac{1}{|\Omega|} \sum_{\omega \in \Omega}
\Big[
\sum_{k=1}^{K}
\big(
C_{W,k}\Gamma_{W,k}(\omega) + C_{L,k}\Gamma_{L,k}(\omega)
\big)
+ \sum_{n=1}^{N}
\big(
C_I^n\Gamma_I^n(\omega) + C_R^n\Gamma_R^n(\omega) + C_O^n\Gamma_O^n(\omega)
\big)
\Big]
\tag{MILP}
\end{equation}

\begin{equation}
\Gamma_{W,k}(\omega) =
\sum_{n=1}^{N} \left( s_k^n(\omega) - r_k^n(\omega) \right),
\qquad \forall k,\ \omega
\tag{10a}
\end{equation}

\begin{equation}
\Gamma_{L,k}(\omega) =
\sum_{n=2}^{N} \left( r_k^n(\omega) - s_k^{n-1}(\omega) - \delta_k^{n-1}(\omega) \right),
\qquad \forall k,\ \omega
\tag{10b}
\end{equation}

\begin{equation}
\Gamma_I^n(\omega) =
\sum_{k=1}^{K} \left( s_k^n(\omega) - v_k^n(\omega) \right),
\qquad \forall n,\ \omega
\tag{10c}
\end{equation}

\begin{equation}
\Gamma_R^n(\omega) =
\sum_{k=2}^{K} \left( v_k^n(\omega) - s_{k-1}^n(\omega) - \delta_{k-1}^n(\omega) \right),
\qquad \forall n,\ \omega
\tag{10d}
\end{equation}

\begin{equation}
\Gamma_O^n(\omega) = f^n(\omega) - T^n,
\qquad \forall n,\ \omega
\tag{10e}
\end{equation}

\begin{equation}
s_k^n(\omega) \ge r_k^n(\omega),
\qquad \forall k,\ n,\ \omega
\tag{11}
\end{equation}

\begin{equation}
s_k^n(\omega) \ge v_k^n(\omega),
\qquad \forall k,\ n,\ \omega
\tag{12}
\end{equation}

\begin{equation}
r_k^1(\omega) = a_k^1,
\qquad \forall k,\ \omega
\tag{13a}
\end{equation}

\begin{equation}
r_k^n(\omega) \ge a_k^n,
\qquad \forall k,\ n \ge 2,\ \omega
\tag{13b}
\end{equation}

\begin{equation}
r_k^n(\omega) \ge s_k^{n-1}(\omega) + \delta_k^{n-1}(\omega),
\qquad \forall k,\ n \ge 2,\ \omega
\tag{14}
\end{equation}

\begin{equation}
r_k^n(\omega) \le a_k^n + M(1-z_k^n(\omega)),
\qquad \forall k,\ n \ge 2,\ \omega
\tag{15}
\end{equation}

\begin{equation}
r_k^n(\omega) \le s_k^{n-1}(\omega) + \delta_k^{n-1}(\omega) + M z_k^n(\omega),
\qquad \forall k,\ n \ge 2,\ \omega
\tag{16}
\end{equation}

\begin{equation}
v_1^n(\omega) = a_1^n,
\qquad \forall n,\ \omega
\tag{17}
\end{equation}

\begin{equation}
v_k^n(\omega) \ge a_k^n,
\qquad \forall k \ge 2,\ n,\ \omega
\tag{18}
\end{equation}

\begin{equation}
v_k^n(\omega) \ge s_{k-1}^n(\omega) + \delta_{k-1}^n(\omega),
\qquad \forall k \ge 2,\ n,\ \omega
\tag{19}
\end{equation}

\begin{equation}
v_k^n(\omega) \le a_k^n + M(1-y_k^n(\omega)),
\qquad \forall k \ge 2,\ n,\ \omega
\tag{20}
\end{equation}

\begin{equation}
v_k^n(\omega) \le s_{k-1}^n(\omega) + \delta_{k-1}^n(\omega) + M y_k^n(\omega),
\qquad \forall k \ge 2,\ n,\ \omega
\tag{21}
\end{equation}

\begin{equation}
f^n(\omega) \ge T^n,
\qquad \forall n,\ \omega
\tag{22}
\end{equation}

\begin{equation}
f^n(\omega) \ge s_K^n(\omega) + \delta_K^n(\omega),
\qquad \forall n,\ \omega
\tag{23}
\end{equation}

\begin{equation}
a_k^n \ge a_{k-1}^n,
\qquad \forall k \ge 2,\ n
\tag{24}
\end{equation}

\begin{equation}
a_k^n \ge a_k^{n-1},
\qquad \forall k,\ n \ge 2
\tag{25}
\end{equation}

\begin{equation}
a_1^1 = 0,\quad
0 \le a_k^n \le T^n,\quad
s_k^n(\omega),\ r_k^n(\omega),\ v_k^n(\omega),\ f^n(\omega) \ge 0,
\qquad \forall k,\ n,\ \omega
\tag{26a}
\end{equation}

\begin{equation}
\Gamma_{W,k}(\omega),\ \Gamma_{L,k}(\omega),\ \Gamma_I^n(\omega),\ \Gamma_R^n(\omega),\ \Gamma_O^n(\omega) \ge 0,
\qquad \forall k,\ n,\ \omega
\tag{26b}
\end{equation}

\begin{equation}
z_k^n(\omega),\ y_k^n(\omega) \in \{0,1\},
\qquad \forall k,\ n,\ \omega
\tag{26c}
\end{equation}


\end{document}
